\begin{lemma}
\label{lemma-make-integral-trick}
Let $R$ be a ring. Let $\varphi : R[x] \to S$ be
a ring map. Let $t \in S$.
Assume that (a) $t$ is integral over $R[x]$,
and (b) there exists a monic $p \in R[x]$ such that
$t \varphi(p) \in \Im(\varphi)$. Then there
exists a $q \in R[x]$ such that $t - \varphi(q)$
is integral over $R$.
\end{lemma}

\begin{proof}
Write $t \varphi(p) = \varphi(r)$ for some $r \in R[x]$.
Using euclidean division, write $r = qp + r'$ with
$q, r' \in R[x]$ and $\deg(r') < \deg(p)$. We may replace
$t$ by $t - \varphi(q)$ which is still integral over
$R[x]$, so that we obtain $t \varphi(p) = \varphi(r')$.
In the ring $S_t$ we may write this as
$\varphi(p) - (1/t) \varphi(r') = 0$.
This implies that $\varphi(x)$ gives an element of the
localization $S_t$ which is integral over
$\varphi(R)[1/t] \subset S_t$. On the other hand,
$t$ is integral over the subring $\varphi(R)[\varphi(x)] \subset S$.
Combined we conclude that $t$ is integral over
the subring $\varphi(R)[1/t] \subset S_t$, see Lemma
\ref{lemma-integral-transitive}. In other words
there exists an equation of the form
$$
t^d + \sum\nolimits_{i < d}
\left(\sum\nolimits_{j = 0, \ldots, n_i} \varphi(r_{i, j})/t^j\right) t^i = 0
$$
in $S_t$ with $r_{i, j} \in R$. This means that
$t^{d + N} +
\sum_{i < d} \sum_{j = 0, \ldots, n_i} \varphi(r_{i, j}) t^{i + N - j} = 0$
in $S$ for some $N$ large enough. In other words
$t$ is integral over $R$.
\end{proof}
